% !TEX root = ../../../main.tex
% Block card: catalogue.minirocket -- "MiniRocket" (Part I, the blocks chapter).
\cardfield{What it does}
Gives every window a shape fingerprint: about ten thousand numbers, each saying how much of
the window a tiny fixed pattern detector fires on.

\begin{plainterms}[A shape fingerprint]
Imagine sliding ten thousand small stencils along a window, each looking for its own little
bump or kink at its own spacing, and writing down for each stencil what fraction of the
window it lit up on. No single number means much. Together they describe the window's shape
well enough that a simple classifier can tell shapes apart from them, quickly
\citep{dempster2021minirocket}.
\end{plainterms}

\cardfield{When to use it}
When you want to know whether a label can be learned from the raw shape of the windows at
all: put it after a windows block and before \blockref{catalogue.classifier}, in place of
(or beside) the hand-picked measures of the window matrix.

\cardfield{What goes in, what comes out}
In goes a window set. Out comes the same window set with the fingerprint added to each
window's row of measures. Windows with too many missing samples are left out and counted.

\cardfield{Settings that matter}
\emph{n\_features} sets how many numbers per window, and with it the time and the table's
size: lower it when a big window set is refused. \emph{normalise} (on) makes the fingerprint
about shape, not level or size. \emph{keep\_upstream\_features} off keeps only the
fingerprint, so a classifier learns from shape alone.

\cardfield{Watch out for}
The detectors' cut-offs are read from training windows only when the window set has a split;
without one they are read from every window and the block says so. A fingerprint with fewer
features that fails to separate your labels is not proof the shapes are the same: the
decisive run is the full 10{,}000 on a shorter span.
